% Cuasiconcavidad: función cóncava, cuasicóncava no cóncava y no cuasicóncava
% (sustituye al escaneo de Pérez Domínguez). En cada panel, la superficie U(x1,x2),
% su corte a la altura k y, en el suelo, el conjunto de contorno superior {x: U(x) >= k}.
% Las curvas de nivel del panel (c) se han calculado numéricamente (U = 0,5).
% Autor: Víctor Gutiérrez Marcos
\documentclass[tikz,border=4pt]{standalone}
% ==========================================================================
% tikz-tcee.tex — Estilos comunes de los gráficos TikZ del temario
% Autor: Víctor Gutiérrez Marcos
%
% Cada gráfico es un fichero standalone en fuentes/<tema>/graficos/:
%
%   \documentclass[tikz,border=4pt]{standalone}
%   % ==========================================================================
% tikz-tcee.tex — Estilos comunes de los gráficos TikZ del temario
% Autor: Víctor Gutiérrez Marcos
%
% Cada gráfico es un fichero standalone en fuentes/<tema>/graficos/:
%
%   \documentclass[tikz,border=4pt]{standalone}
%   % ==========================================================================
% tikz-tcee.tex — Estilos comunes de los gráficos TikZ del temario
% Autor: Víctor Gutiérrez Marcos
%
% Cada gráfico es un fichero standalone en fuentes/<tema>/graficos/:
%
%   \documentclass[tikz,border=4pt]{standalone}
%   \input{../../../latex/tikz-tcee}
%   \begin{document}
%   \begin{tikzpicture}[tcee]
%     \ejes{Q}{P}{6}{5}
%     \draw[curva1] (0.5,4.5) -- (5,0.8) node[etiqueta,right] {$D$};
%     \capa{2}{\draw[curva2] ...;}   % en el vídeo aparece en el paso 2
%   \end{tikzpicture}
%   \end{document}
%
% Paleta: la de la web (morado, dorado, salvia) más tres colores de apoyo
% distinguibles en blanco y negro por el trazo.
% ==========================================================================
\usepackage[T1]{fontenc}
\usepackage{newpxtext,newpxmath}
\usepackage{xcolor}
\usepackage{tikz,pgfplots}
\pgfplotsset{compat=1.18}
\usetikzlibrary{arrows.meta,calc,positioning,patterns,decorations.pathreplacing,intersections,shapes.misc,backgrounds}

\definecolor{tceemorado}{HTML}{5F2987}
\definecolor{tceedorado}{HTML}{B8860B}
\definecolor{tceeverde}{HTML}{3C7D3A}
\definecolor{tceeazul}{HTML}{1F5F99}
\definecolor{tceerojo}{HTML}{B22222}
\definecolor{tceegris}{HTML}{767171}
\definecolor{tceesalvia}{HTML}{E2EFD9}

% Capas para el vídeo: \capa{n}{…} dibuja su contenido solo a partir del paso n.
% En el PDF, la web y Word se ve todo; generar-video.py compila el gráfico una
% vez por paso con \def\tceepaso{k} para que cada elemento aparezca al narrarlo.
\providecommand{\tceepaso}{1000}
\newcommand{\capa}[2]{\ifnum#1>\tceepaso\relax\else#2\fi}

\tikzset{
  tcee/.style={line cap=round,line join=round,font=\small,>={Stealth[length=2.2mm]}},
  eje/.style={->,thick,black},
  curva1/.style={very thick,tceemorado},
  curva2/.style={very thick,tceeazul},
  curva3/.style={very thick,tceeverde},
  curva4/.style={very thick,tceedorado},
  curvanueva/.style={very thick,tceerojo},
  desplazada/.style={very thick,dashed},
  guia/.style={thin,dashed,tceegris},
  punto/.style={circle,fill=black,inner sep=1.3pt},
  etiqueta/.style={font=\small,inner sep=2pt},
  flecha/.style={->,thick,tceerojo},
  area/.style={fill=tceemorado,fill opacity=0.18,draw=none},
  area2/.style={fill=tceedorado,fill opacity=0.22,draw=none},
}

% \ejes{etiqueta x}{etiqueta y}{largo x}{alto y}
\newcommand{\ejes}[4]{%
  \draw[eje] (0,0) -- (#3,0) node[below left=2pt and 0pt] {#1};
  \draw[eje] (0,0) -- (0,#4) node[left] {#2};
  \node[below left] at (0,0) {$0$};}
% \proyeccion{(x,y)}{etq x}{etq y}: líneas guía a los ejes con etiquetas
\newcommand{\proyeccion}[3]{%
  \draw[guia] let \p1=#1 in (\x1,0) node[below] {#2} |- (0,\y1) node[left] {#3};}

%   \begin{document}
%   \begin{tikzpicture}[tcee]
%     \ejes{Q}{P}{6}{5}
%     \draw[curva1] (0.5,4.5) -- (5,0.8) node[etiqueta,right] {$D$};
%     \capa{2}{\draw[curva2] ...;}   % en el vídeo aparece en el paso 2
%   \end{tikzpicture}
%   \end{document}
%
% Paleta: la de la web (morado, dorado, salvia) más tres colores de apoyo
% distinguibles en blanco y negro por el trazo.
% ==========================================================================
\usepackage[T1]{fontenc}
\usepackage{newpxtext,newpxmath}
\usepackage{xcolor}
\usepackage{tikz,pgfplots}
\pgfplotsset{compat=1.18}
\usetikzlibrary{arrows.meta,calc,positioning,patterns,decorations.pathreplacing,intersections,shapes.misc,backgrounds}

\definecolor{tceemorado}{HTML}{5F2987}
\definecolor{tceedorado}{HTML}{B8860B}
\definecolor{tceeverde}{HTML}{3C7D3A}
\definecolor{tceeazul}{HTML}{1F5F99}
\definecolor{tceerojo}{HTML}{B22222}
\definecolor{tceegris}{HTML}{767171}
\definecolor{tceesalvia}{HTML}{E2EFD9}

% Capas para el vídeo: \capa{n}{…} dibuja su contenido solo a partir del paso n.
% En el PDF, la web y Word se ve todo; generar-video.py compila el gráfico una
% vez por paso con \def\tceepaso{k} para que cada elemento aparezca al narrarlo.
\providecommand{\tceepaso}{1000}
\newcommand{\capa}[2]{\ifnum#1>\tceepaso\relax\else#2\fi}

\tikzset{
  tcee/.style={line cap=round,line join=round,font=\small,>={Stealth[length=2.2mm]}},
  eje/.style={->,thick,black},
  curva1/.style={very thick,tceemorado},
  curva2/.style={very thick,tceeazul},
  curva3/.style={very thick,tceeverde},
  curva4/.style={very thick,tceedorado},
  curvanueva/.style={very thick,tceerojo},
  desplazada/.style={very thick,dashed},
  guia/.style={thin,dashed,tceegris},
  punto/.style={circle,fill=black,inner sep=1.3pt},
  etiqueta/.style={font=\small,inner sep=2pt},
  flecha/.style={->,thick,tceerojo},
  area/.style={fill=tceemorado,fill opacity=0.18,draw=none},
  area2/.style={fill=tceedorado,fill opacity=0.22,draw=none},
}

% \ejes{etiqueta x}{etiqueta y}{largo x}{alto y}
\newcommand{\ejes}[4]{%
  \draw[eje] (0,0) -- (#3,0) node[below left=2pt and 0pt] {#1};
  \draw[eje] (0,0) -- (0,#4) node[left] {#2};
  \node[below left] at (0,0) {$0$};}
% \proyeccion{(x,y)}{etq x}{etq y}: líneas guía a los ejes con etiquetas
\newcommand{\proyeccion}[3]{%
  \draw[guia] let \p1=#1 in (\x1,0) node[below] {#2} |- (0,\y1) node[left] {#3};}

%   \begin{document}
%   \begin{tikzpicture}[tcee]
%     \ejes{Q}{P}{6}{5}
%     \draw[curva1] (0.5,4.5) -- (5,0.8) node[etiqueta,right] {$D$};
%     \capa{2}{\draw[curva2] ...;}   % en el vídeo aparece en el paso 2
%   \end{tikzpicture}
%   \end{document}
%
% Paleta: la de la web (morado, dorado, salvia) más tres colores de apoyo
% distinguibles en blanco y negro por el trazo.
% ==========================================================================
\usepackage[T1]{fontenc}
\usepackage{newpxtext,newpxmath}
\usepackage{xcolor}
\usepackage{tikz,pgfplots}
\pgfplotsset{compat=1.18}
\usetikzlibrary{arrows.meta,calc,positioning,patterns,decorations.pathreplacing,intersections,shapes.misc,backgrounds}

\definecolor{tceemorado}{HTML}{5F2987}
\definecolor{tceedorado}{HTML}{B8860B}
\definecolor{tceeverde}{HTML}{3C7D3A}
\definecolor{tceeazul}{HTML}{1F5F99}
\definecolor{tceerojo}{HTML}{B22222}
\definecolor{tceegris}{HTML}{767171}
\definecolor{tceesalvia}{HTML}{E2EFD9}

% Capas para el vídeo: \capa{n}{…} dibuja su contenido solo a partir del paso n.
% En el PDF, la web y Word se ve todo; generar-video.py compila el gráfico una
% vez por paso con \def\tceepaso{k} para que cada elemento aparezca al narrarlo.
\providecommand{\tceepaso}{1000}
\newcommand{\capa}[2]{\ifnum#1>\tceepaso\relax\else#2\fi}

\tikzset{
  tcee/.style={line cap=round,line join=round,font=\small,>={Stealth[length=2.2mm]}},
  eje/.style={->,thick,black},
  curva1/.style={very thick,tceemorado},
  curva2/.style={very thick,tceeazul},
  curva3/.style={very thick,tceeverde},
  curva4/.style={very thick,tceedorado},
  curvanueva/.style={very thick,tceerojo},
  desplazada/.style={very thick,dashed},
  guia/.style={thin,dashed,tceegris},
  punto/.style={circle,fill=black,inner sep=1.3pt},
  etiqueta/.style={font=\small,inner sep=2pt},
  flecha/.style={->,thick,tceerojo},
  area/.style={fill=tceemorado,fill opacity=0.18,draw=none},
  area2/.style={fill=tceedorado,fill opacity=0.22,draw=none},
}

% \ejes{etiqueta x}{etiqueta y}{largo x}{alto y}
\newcommand{\ejes}[4]{%
  \draw[eje] (0,0) -- (#3,0) node[below left=2pt and 0pt] {#1};
  \draw[eje] (0,0) -- (0,#4) node[left] {#2};
  \node[below left] at (0,0) {$0$};}
% \proyeccion{(x,y)}{etq x}{etq y}: líneas guía a los ejes con etiquetas
\newcommand{\proyeccion}[3]{%
  \draw[guia] let \p1=#1 in (\x1,0) node[below] {#2} |- (0,\y1) node[left] {#3};}

\pgfplotsset{
  panel/.style={width=6.4cm,height=6.4cm,view/h=30,view/v=28,axis lines=none,
    zmin=-0.9,clip=false,title style={yshift=-1.3cm},
    colormap={tcee}{rgb255=(226,239,217) rgb255=(167,137,189)}},
  superficie/.style={surf,shader=faceted,fill opacity=0.8,draw=tceemorado!45,line width=0.15pt},
  suelo/.style={fill=tceemorado,fill opacity=0.3,draw=tceemorado,thick},
  corte/.style={curvanueva,line width=1pt},
}
\newcommand{\ejespanel}[3]{% #1 = x max, #2 = y max, #3 = z max
  \draw[eje] (axis cs:-#1,-#2,-0.9) -- (axis cs:#1+0.5,-#2,-0.9) node[below right=-2pt] {$x_1$};
  \draw[eje] (axis cs:-#1,-#2,-0.9) -- (axis cs:-#1,#2+0.6,-0.9);
  \draw[eje] (axis cs:-#1,-#2,-0.9) -- (axis cs:-#1,-#2,#3) node[above] {$U$};}
\begin{document}
\begin{tikzpicture}[tcee]
  % ---------- 1: función cóncava ----------
  \capa{1}{
  \begin{axis}[panel,at={(0,0)},title={(a) Cóncava},zmax=1.15]
    \ejespanel{1.8}{1.8}{1.25}
    \addplot3[suelo] coordinates {(1.200,0.000,-0.9) (1.182,0.208,-0.9) (1.128,0.410,-0.9) (1.039,0.600,-0.9) (0.919,0.771,-0.9) (0.771,0.919,-0.9) (0.600,1.039,-0.9) (0.410,1.128,-0.9) (0.208,1.182,-0.9) (0.000,1.200,-0.9) (-0.208,1.182,-0.9) (-0.410,1.128,-0.9) (-0.600,1.039,-0.9) (-0.771,0.919,-0.9) (-0.919,0.771,-0.9) (-1.039,0.600,-0.9) (-1.128,0.410,-0.9) (-1.182,0.208,-0.9) (-1.200,0.000,-0.9) (-1.182,-0.208,-0.9) (-1.128,-0.410,-0.9) (-1.039,-0.600,-0.9) (-0.919,-0.771,-0.9) (-0.771,-0.919,-0.9) (-0.600,-1.039,-0.9) (-0.410,-1.128,-0.9) (-0.208,-1.182,-0.9) (-0.000,-1.200,-0.9) (0.208,-1.182,-0.9) (0.410,-1.128,-0.9) (0.600,-1.039,-0.9) (0.771,-0.919,-0.9) (0.919,-0.771,-0.9) (1.039,-0.600,-0.9) (1.128,-0.410,-0.9) (1.182,-0.208,-0.9) (1.200,-0.000,-0.9)};
    \addplot3[superficie,domain=0:1.8,y domain=0:360,samples=13,samples y=29]
      ({x*cos(y)},{x*sin(y)},{1-x^2/3.6});
    \addplot3[corte,domain=0:360,samples=60] ({1.2*cos(x)},{1.2*sin(x)},0.6);
    \node[etiqueta,text=tceerojo,anchor=west] at (axis cs:1.9,0,0.6) {$U=k$};
    \node[etiqueta,right] at (axis cs:-1.8,2.4,-0.9) {$x_2$};
  \end{axis}}
  % ---------- 2: función cuasicóncava no cóncava ----------
  \capa{2}{
  \begin{axis}[panel,at={(6.2cm,0)},title={(b) Cuasicóncava, no cóncava},zmax=1.15]
    \ejespanel{2.2}{2.2}{1.25}
    \addplot3[suelo] coordinates {(0.957,0.000,-0.9) (0.942,0.166,-0.9) (0.899,0.327,-0.9) (0.829,0.478,-0.9) (0.733,0.615,-0.9) (0.615,0.733,-0.9) (0.479,0.829,-0.9) (0.327,0.899,-0.9) (0.166,0.942,-0.9) (0.000,0.957,-0.9) (-0.166,0.942,-0.9) (-0.327,0.899,-0.9) (-0.478,0.829,-0.9) (-0.615,0.733,-0.9) (-0.733,0.615,-0.9) (-0.829,0.478,-0.9) (-0.899,0.327,-0.9) (-0.942,0.166,-0.9) (-0.957,0.000,-0.9) (-0.942,-0.166,-0.9) (-0.899,-0.327,-0.9) (-0.829,-0.479,-0.9) (-0.733,-0.615,-0.9) (-0.615,-0.733,-0.9) (-0.479,-0.829,-0.9) (-0.327,-0.899,-0.9) (-0.166,-0.942,-0.9) (-0.000,-0.957,-0.9) (0.166,-0.942,-0.9) (0.327,-0.899,-0.9) (0.479,-0.829,-0.9) (0.615,-0.733,-0.9) (0.733,-0.615,-0.9) (0.829,-0.479,-0.9) (0.899,-0.327,-0.9) (0.942,-0.166,-0.9) (0.957,-0.000,-0.9)};
    \addplot3[superficie,domain=0:2.2,y domain=0:360,samples=15,samples y=29]
      ({x*cos(y)},{x*sin(y)},{exp(-x^2)});
    \addplot3[corte,domain=0:360,samples=60] ({0.957*cos(x)},{0.957*sin(x)},0.4);
    \node[etiqueta,text=tceerojo,anchor=west] at (axis cs:2.3,0,0.4) {$U=k$};
    \node[etiqueta,right] at (axis cs:-2.2,2.8,-0.9) {$x_2$};
  \end{axis}}
  % ---------- 3: función no cuasicóncava ----------
  \capa{3}{
  \begin{axis}[panel,at={(12.4cm,0)},title={(c) No cuasicóncava},zmax=1.15]
    \ejespanel{3}{2}{1.25}
    \addplot3[suelo] coordinates {(2.133,0.000,-0.9) (2.120,0.145,-0.9) (2.082,0.285,-0.9) (2.021,0.416,-0.9) (1.938,0.535,-0.9) (1.835,0.638,-0.9) (1.716,0.721,-0.9) (1.585,0.782,-0.9) (1.445,0.820,-0.9) (1.300,0.833,-0.9) (1.155,0.821,-0.9) (1.014,0.785,-0.9) (0.881,0.726,-0.9) (0.757,0.647,-0.9) (0.647,0.548,-0.9) (0.551,0.433,-0.9) (0.473,0.301,-0.9) (0.419,0.155,-0.9) (0.399,0.000,-0.9) (0.419,-0.155,-0.9) (0.473,-0.301,-0.9) (0.551,-0.433,-0.9) (0.647,-0.548,-0.9) (0.757,-0.647,-0.9) (0.881,-0.726,-0.9) (1.014,-0.785,-0.9) (1.155,-0.821,-0.9) (1.300,-0.833,-0.9) (1.445,-0.820,-0.9) (1.585,-0.782,-0.9) (1.716,-0.721,-0.9) (1.835,-0.638,-0.9) (1.938,-0.535,-0.9) (2.021,-0.416,-0.9) (2.082,-0.285,-0.9) (2.120,-0.145,-0.9) (2.133,0.000,-0.9)};
    \addplot3[suelo] coordinates {(-0.399,0.000,-0.9) (-0.419,0.155,-0.9) (-0.473,0.301,-0.9) (-0.551,0.433,-0.9) (-0.647,0.548,-0.9) (-0.757,0.647,-0.9) (-0.881,0.726,-0.9) (-1.014,0.785,-0.9) (-1.155,0.821,-0.9) (-1.300,0.833,-0.9) (-1.445,0.820,-0.9) (-1.585,0.782,-0.9) (-1.716,0.721,-0.9) (-1.835,0.638,-0.9) (-1.938,0.535,-0.9) (-2.021,0.416,-0.9) (-2.082,0.285,-0.9) (-2.120,0.145,-0.9) (-2.133,0.000,-0.9) (-2.120,-0.145,-0.9) (-2.082,-0.285,-0.9) (-2.021,-0.416,-0.9) (-1.938,-0.535,-0.9) (-1.835,-0.638,-0.9) (-1.716,-0.721,-0.9) (-1.585,-0.782,-0.9) (-1.445,-0.820,-0.9) (-1.300,-0.833,-0.9) (-1.155,-0.821,-0.9) (-1.014,-0.785,-0.9) (-0.881,-0.726,-0.9) (-0.757,-0.647,-0.9) (-0.647,-0.548,-0.9) (-0.551,-0.433,-0.9) (-0.473,-0.301,-0.9) (-0.419,-0.155,-0.9) (-0.399,0.000,-0.9)};
    \addplot3[superficie,domain=-3:3,y domain=-2:2,samples=25,samples y=17]
      {exp(-((x-1.3)^2+y^2))+exp(-((x+1.3)^2+y^2))};
    \addplot3[corte] coordinates {(2.133,0.000,0.5) (2.120,0.145,0.5) (2.082,0.285,0.5) (2.021,0.416,0.5) (1.938,0.535,0.5) (1.835,0.638,0.5) (1.716,0.721,0.5) (1.585,0.782,0.5) (1.445,0.820,0.5) (1.300,0.833,0.5) (1.155,0.821,0.5) (1.014,0.785,0.5) (0.881,0.726,0.5) (0.757,0.647,0.5) (0.647,0.548,0.5) (0.551,0.433,0.5) (0.473,0.301,0.5) (0.419,0.155,0.5) (0.399,0.000,0.5) (0.419,-0.155,0.5) (0.473,-0.301,0.5) (0.551,-0.433,0.5) (0.647,-0.548,0.5) (0.757,-0.647,0.5) (0.881,-0.726,0.5) (1.014,-0.785,0.5) (1.155,-0.821,0.5) (1.300,-0.833,0.5) (1.445,-0.820,0.5) (1.585,-0.782,0.5) (1.716,-0.721,0.5) (1.835,-0.638,0.5) (1.938,-0.535,0.5) (2.021,-0.416,0.5) (2.082,-0.285,0.5) (2.120,-0.145,0.5) (2.133,0.000,0.5)};
    \addplot3[corte] coordinates {(-0.399,0.000,0.5) (-0.419,0.155,0.5) (-0.473,0.301,0.5) (-0.551,0.433,0.5) (-0.647,0.548,0.5) (-0.757,0.647,0.5) (-0.881,0.726,0.5) (-1.014,0.785,0.5) (-1.155,0.821,0.5) (-1.300,0.833,0.5) (-1.445,0.820,0.5) (-1.585,0.782,0.5) (-1.716,0.721,0.5) (-1.835,0.638,0.5) (-1.938,0.535,0.5) (-2.021,0.416,0.5) (-2.082,0.285,0.5) (-2.120,0.145,0.5) (-2.133,0.000,0.5) (-2.120,-0.145,0.5) (-2.082,-0.285,0.5) (-2.021,-0.416,0.5) (-1.938,-0.535,0.5) (-1.835,-0.638,0.5) (-1.716,-0.721,0.5) (-1.585,-0.782,0.5) (-1.445,-0.820,0.5) (-1.300,-0.833,0.5) (-1.155,-0.821,0.5) (-1.014,-0.785,0.5) (-0.881,-0.726,0.5) (-0.757,-0.647,0.5) (-0.647,-0.548,0.5) (-0.551,-0.433,0.5) (-0.473,-0.301,0.5) (-0.419,-0.155,0.5) (-0.399,0.000,0.5)};
    \node[etiqueta,text=tceerojo,anchor=west] at (axis cs:3.1,0,0.5) {$U=k$};
    \node[etiqueta,right] at (axis cs:-3,2.6,-0.9) {$x_2$};
  \end{axis}}
  % rótulo común del suelo
  \node[etiqueta,align=center,font=\footnotesize] at (9.3cm,-1.0cm)
    {En el suelo, sombreado, el conjunto de contorno superior $\{x: U(x)\geq k\}$: convexo en (a) y (b), no convexo en (c)};
\end{tikzpicture}
\end{document}
